%! The Nullspace of A: Solving Ax = 0 — Unit 3, Lesson 3.2 — Exit Ticket (Answer Key)
\documentclass[10pt]{article}
\usepackage{linalg-article}
\usepackage{linalg-key}   % pulls in -boxes; do NOT also load -boxes
\newcommand{\TallMath}[1]{$\displaystyle #1\rule[-1.4em]{0pt}{3.2em}$}
\newcommand{\xx}{\mathbf{x}}
\newcommand{\bb}{\mathbf{b}}
\newcommand{\svec}{\mathbf{s}}
\newcommand{\zero}{\mathbf{0}}

\begin{document}

\pageheader{Unit 3, Lesson 3.2}{Exit Ticket --- Answer Key}
\namedateperiod

\vspace{0.15in}

No notes. Show your reasoning.

\begin{enumerate}[leftmargin=1.8em, itemsep=16pt]
  \item \textbf{Find the nullspace.} For $A=\begin{bmatrix} 1 & 4 \\ 2 & 8 \end{bmatrix}$, reduce $A$,
        give the special solution, and say whether $N(A)$ is a point, line, or plane.
        \ansline{$R_2-2R_1\Rightarrow\begin{bsmallmatrix} 1&4\\0&0 \end{bsmallmatrix}$, so $x_1+4x_2=0$ with $x_2$ free. Special solution $\svec=\begin{bsmallmatrix} -4\\1 \end{bsmallmatrix}$; $N(A)$ = all multiples of $\svec$ = a \emph{line} through the origin ($n-r=2-1=1$).}
  \item \textbf{Count the free variables.} A matrix $A$ has $4$ columns and rank $r=2$. How many free
        variables does $A\xx=\zero$ have, and what is the dimension of $N(A)$?
        \ansline{Free variables $=n-r=4-2=2$; so $N(A)$ has $2$ special solutions and dimension $2$ (a plane through the origin).}
  \item \textbf{Justify.} A $3\times3$ matrix $A$ has nullspace $N(A)=\{\zero\}$ (only the zero vector).
        What does this tell you about how many solutions $A\xx=\bb$ has, and whether $A$ is invertible?
        Explain.
        \ansline{With no free variables ($r=n=3$), the only solution of $A\xx=\zero$ is $\zero$, so $A\xx=\bb$ has \emph{exactly one} solution for every $\bb$ --- $A$ is \emph{invertible}. (Two different solutions would differ by a nonzero nullspace vector, which does not exist here.)}
\end{enumerate}

\begin{teachernote}
Sort responses into ``got it / free-vs-pivot slip / stopped at one special solution.'' Item~3 is the
discriminator and the bridge to 3.3: students should connect ``$N(A)=\{\zero\}$'' to ``unique
solution / invertible,'' not just restate the definition. If many stumble on item~1's reduction, open
3.3 by re-finding a nullspace before adding a right-hand side.
\end{teachernote}

\end{document}
